%% !TEX program = lualatex
% ============================================================================
% Human in the Loop: Misunderstood — Spring Edition
% Semester Project Guide
% ============================================================================
% Compile: lualatex project_guide.tex (twice for TOC/refs)
% ============================================================================

\documentclass[
  11pt,
  oneside,
  DIV=12,
]{scrartcl}

% ── Fonts ────────────────────────────────────────────────────────────────────
\usepackage{fontspec}
\usepackage{microtype}
\setmainfont[
  Path=/usr/local/texlive/2025/texmf-dist/fonts/opentype/public/libertinus-fonts/,
  BoldFont=LibertinusSerif-Bold.otf,
  ItalicFont=LibertinusSerif-Italic.otf,
  BoldItalicFont=LibertinusSerif-BoldItalic.otf,
  Numbers=OldStyle,
  Ligatures=TeX,
]{LibertinusSerif-Regular.otf}
\setsansfont[
  Path=/usr/local/texlive/2025/texmf-dist/fonts/opentype/public/libertinus-fonts/,
  BoldFont=LibertinusSans-Bold.otf,
  ItalicFont=LibertinusSans-Italic.otf,
]{LibertinusSans-Regular.otf}
\setmonofont[
  Path=/usr/local/texlive/2025/texmf-dist/fonts/opentype/SIL/jetbrainsmono-otf/,
  BoldFont=JetBrainsMono-Bold.otf,
  ItalicFont=JetBrainsMono-Italic.otf,
  Scale=0.82,
  Contextuals=Alternate,
]{JetBrainsMono-Regular.otf}

% ── Geometry ─────────────────────────────────────────────────────────────────
\usepackage{geometry}
\geometry{
  paper=a4paper,
  top=2.4cm,
  bottom=2.8cm,
  inner=2.8cm,
  outer=2.8cm,
}

% ── Colors ───────────────────────────────────────────────────────────────────
\usepackage{xcolor}
\definecolor{SpringGreen}{HTML}{2D6A4F}
\definecolor{SpringLeaf}{HTML}{40916C}
\definecolor{SpringFog}{HTML}{D8F3DC}
\definecolor{SpringMoss}{HTML}{1B4332}
\definecolor{HintAmber}{HTML}{B45309}
\definecolor{HintCream}{HTML}{FEF9EF}
\definecolor{WarnRed}{HTML}{9B1C1C}
\definecolor{WarnPink}{HTML}{FEE2E2}
\definecolor{ConceptBlue}{HTML}{1E3A5F}
\definecolor{ConceptIce}{HTML}{EBF4FF}
\definecolor{DeliverGray}{HTML}{374151}
\definecolor{DeliverFog}{HTML}{F3F4F6}
\definecolor{RuleGray}{HTML}{D1D5DB}
\definecolor{WeekBand}{HTML}{F9FAFB}

% ── Packages ─────────────────────────────────────────────────────────────────
\usepackage{tcolorbox}
\tcbuselibrary{skins, breakable, listings}
\usepackage{enumitem}
\usepackage{booktabs}
\usepackage{array}
\usepackage{tabularx}
\usepackage{multicol}
\usepackage{hyperref}
\usepackage{polyglossia}
\setdefaultlanguage[variant=american]{english}
\usepackage{fontawesome5}

\hypersetup{
  colorlinks=true,
  linkcolor=SpringMoss,
  urlcolor=SpringLeaf,
  pdftitle={Human in the Loop: Misunderstood — Spring Edition: Semester Project Guide},
  pdfauthor={Lazaros Toumanidis},
}

% ── Typography ───────────────────────────────────────────────────────────────
\setlength{\parindent}{0pt}
\setlength{\parskip}{6pt}
\setlist[itemize]{leftmargin=1.4em, topsep=3pt, itemsep=1pt}
\setlist[enumerate]{leftmargin=1.6em, topsep=3pt, itemsep=1pt}

% ── tcolorbox environments ───────────────────────────────────────────────────

% Milestone card — the main building block
\newtcolorbox{milestone}[2]{
  enhanced,
  breakable,
  before skip=10pt, after skip=6pt,
  colframe=SpringLeaf,
  colback=white,
  colbacktitle=SpringGreen,
  coltitle=white,
  fonttitle=\sffamily\bfseries\large,
  title={\faLeaf\quad Week~#1 \enspace\textbar\enspace #2},
  top=8pt, bottom=8pt, left=10pt, right=10pt,
  arc=4pt,
  drop shadow,
}

% HITL concept callout inside milestone
\newtcolorbox{hitlconcept}[1]{
  enhanced,
  colframe=ConceptBlue!40,
  colback=ConceptIce,
  coltitle=ConceptBlue,
  fonttitle=\sffamily\bfseries\small,
  title={\faLightbulb[regular]\enspace HITL Concept: #1},
  top=4pt, bottom=4pt, left=8pt, right=8pt,
  arc=3pt,
  boxrule=0.6pt,
  before skip=5pt, after skip=3pt,
}

% Hint box
\newtcolorbox{hintbox}{
  enhanced,
  colframe=HintAmber!60,
  colback=HintCream,
  coltitle=HintAmber,
  fonttitle=\sffamily\bfseries\small,
  title={\faQuestion\enspace Hint},
  top=3pt, bottom=3pt, left=8pt, right=8pt,
  arc=3pt,
  boxrule=0.5pt,
  before skip=4pt, after skip=2pt,
}

% Checkpoint deliverable
\newtcolorbox{checkpoint}{
  enhanced,
  colframe=DeliverGray!40,
  colback=DeliverFog,
  coltitle=DeliverGray,
  fonttitle=\sffamily\bfseries\small,
  title={\faCheckCircle[regular]\enspace Checkpoint},
  top=3pt, bottom=3pt, left=8pt, right=8pt,
  arc=3pt,
  boxrule=0.5pt,
  before skip=4pt, after skip=2pt,
}

% Caution / watch-out box
\newtcolorbox{watchout}{
  enhanced,
  colframe=WarnRed!50,
  colback=WarnPink,
  coltitle=WarnRed,
  fonttitle=\sffamily\bfseries\small,
  title={\faExclamationTriangle\enspace Watch out},
  top=3pt, bottom=3pt, left=8pt, right=8pt,
  arc=3pt,
  boxrule=0.5pt,
  before skip=4pt, after skip=2pt,
}

% Overview/info box
\newtcolorbox{infobox}[1]{
  enhanced,
  colframe=SpringLeaf!50,
  colback=SpringFog!40,
  coltitle=SpringMoss,
  fonttitle=\sffamily\bfseries,
  title={#1},
  top=6pt, bottom=6pt, left=10pt, right=10pt,
  arc=4pt,
  boxrule=0.8pt,
  before skip=8pt, after skip=8pt,
}

% ── Section styling ───────────────────────────────────────────────────────────
\RedeclareSectionCommand[
  beforeskip=24pt,
  afterskip=6pt,
  font=\Large\sffamily\bfseries\color{SpringMoss},
]{section}
\RedeclareSectionCommand[
  beforeskip=14pt,
  afterskip=4pt,
  font=\large\sffamily\bfseries\color{SpringGreen},
]{subsection}

% ── Helpers ───────────────────────────────────────────────────────────────────
\newcommand{\build}[1]{\textbf{\color{SpringMoss}Build:} #1}
\newcommand{\reflect}[1]{\textbf{\color{ConceptBlue}Reflect:} \textit{#1}}
\newcommand{\role}[1]{\textsc{\color{SpringLeaf}#1}}

% ── Title block ───────────────────────────────────────────────────────────────
\setkomafont{title}{\Huge\sffamily\bfseries\color{SpringMoss}}
\setkomafont{subtitle}{\large\sffamily\color{SpringLeaf}}
\setkomafont{author}{\normalsize\sffamily}
\setkomafont{date}{\small\sffamily\color{gray}}

% ============================================================================
\begin{document}
% ============================================================================

% ── Cover ─────────────────────────────────────────────────────────────────────
\thispagestyle{empty}
\begin{center}
  \vspace*{1.2cm}

  {\color{RuleGray}\rule{\linewidth}{0.4pt}}
  \vspace{0.5cm}

  {\fontsize{32}{38}\selectfont\sffamily\bfseries\color{SpringMoss}
    Build the Loop}

  \vspace{0.3cm}

  {\LARGE\sffamily\color{SpringLeaf}
    Spring Edition — Semester Project Guide}

  \vspace{0.4cm}

  {\color{RuleGray}\rule{\linewidth}{0.4pt}}

  \vspace{0.8cm}

  {\large\sffamily\color{DeliverGray}
    Human in the Loop: Misunderstood}

  \vspace{0.3cm}

  {\normalsize\sffamily\color{gray}
    Lazaros Toumanidis}

  \vspace{1.2cm}

  \begin{infobox}{In one sentence}
    Over fifteen weeks, you will build a complete Human-in-the-Loop system
    from a blank directory to a live, deployed treasure hunt — and in doing
    so, you will \emph{feel} every concept that the other editions only
    describe.
  \end{infobox}

  \vspace{0.8cm}

  {\small\sffamily\color{gray}
    \faTools\enspace Rust + Axum \quad
    \faPython\enspace Python + FastAPI \quad
    \faReact\enspace React \quad
    \faMobile*\enspace Flutter \quad
    \faDocker\enspace Docker}

\end{center}

\clearpage

% ── Overview ─────────────────────────────────────────────────────────────────
\section*{What This Project Is}
\addcontentsline{toc}{section}{What This Project Is}

The Spring Edition is different from every other edition of this series.

The Summer Edition asks you to \emph{think}. The Winter Edition asks you to
\emph{analyse}. This one asks you to \emph{build} — and then to run what you
built with real people in a real space, and watch what breaks.

The system is a \textbf{Treasure Hunt}: players follow a sequence of clues
through a physical location, scanning QR codes or tapping NFC tags, submitting
answers, and earning hints when they are stuck. Behind the game is a complete
HITL architecture: a \role{creator} who designs the experience, a \role{player}
who generates feedback by playing, and a \role{developer} (you) who builds the
machinery connecting them through an AI layer.

\begin{infobox}{Why a game?}
  Games have four properties that make HITL legible:
  \begin{itemize}
    \item \textbf{Defined tasks} — a clue has a correct answer, so we can
          measure human performance without ambiguity.
    \item \textbf{Motivated participants} — players want to win; their
          behavior is signal, not noise.
    \item \textbf{Visible feedback} — attempts, timing, and hint requests
          are explicit training signal.
    \item \textbf{Low stakes} — wrong predictions mean nudges, not harm.
          We can iterate freely and break things deliberately.
  \end{itemize}
\end{infobox}

\subsection*{The Three Roles}

\begin{tabularx}{\linewidth}{@{}>{\bfseries\color{SpringMoss}}lX@{}}
\toprule
\role{Creator} & The annotator. Encodes judgment into data by designing clues,
  setting thresholds, choosing what counts as correct.
  Their decisions are baked in long before any player sees them. \\[4pt]
\role{Player}  & The oracle. Receives outputs (clues) and produces feedback
  (attempts, timing, hints). Does not know they are providing training signal.
  They are just trying to win. \\[4pt]
\role{Developer} & You. Builds the infrastructure connecting creator and
  player through an AI layer, and must decide how much autonomy to give
  each component. \\
\bottomrule
\end{tabularx}

\bigskip

\subsection*{What You Will Deliver}

By the end of week 15, you will have:

\begin{itemize}
  \item A \textbf{Rust/Axum} REST API with WebSocket live events
  \item A \textbf{Python/FastAPI} AI sidecar (clue generation, difficulty
        estimation, adaptive hints)
  \item A \textbf{React} creator dashboard
  \item A \textbf{Flutter} mobile app for players
  \item A running local stack via
        \texttt{docker compose up}
  \item A 2-page \textbf{retrospective analysis} mapping what you built
        to the Five Dimensions
\end{itemize}

\clearpage

% ── Prerequisites ─────────────────────────────────────────────────────────────
\section*{Prerequisites and Setup}

\begin{infobox}{Required before Week 1}
\begin{itemize}
  \item Rust toolchain (\texttt{rustup}, stable channel)
  \item Python 3.11+
  \item Docker + Docker Compose v2
  \item Node.js 20+ and npm
  \item Flutter SDK 3.22+ (Chapters 10, 12 only)
  \item An LLM API key if you want hosted-model features (\texttt{LLM\_API\_KEY})
\end{itemize}
\end{infobox}

\subsection*{Getting Started}

\begin{enumerate}
  \item Clone the repository and enter the Spring Edition directory:
  \begin{center}\ttfamily\small
    cd path/to/repo/editions/spring
  \end{center}

  \item Copy the environment template and fill in your model settings:
  \begin{center}\ttfamily\small
    cp .env.example .env
  \end{center}

  \item Start the local stack:
  \begin{center}\ttfamily\small
    docker compose up -d --build
  \end{center}

  \item Verify Rust compiles (takes 2--5 minutes on first build):
  \begin{center}\ttfamily\small
    cd backend-rust \&\& cargo build
  \end{center}

  \item Install Python dependencies:
  \begin{center}\ttfamily\small
    cd ../backend-python \&\& pip install -r requirements.txt
  \end{center}
\end{enumerate}

If \texttt{cargo build} succeeds in your environment and \texttt{docker compose ps} shows
the services running, you are ready.

\begin{watchout}
  The Rust build downloads dependencies on first run.
  Do this before Week 1 class, not during it.
\end{watchout}

\clearpage

% ── 15-Week Schedule ─────────────────────────────────────────────────────────
\section*{15-Week Schedule at a Glance}

\begin{center}
\renewcommand{\arraystretch}{1.35}
\begin{tabularx}{\linewidth}{@{}>{\sffamily\bfseries\color{SpringMoss}}lX>{\small\color{SpringLeaf}\sffamily}l@{}}
\toprule
Week & Chapter & HITL Concept \\
\midrule
1  & Building HITL: Why a Game? & The Loop Made Tangible \\
2  & The Five Dimensions in Play & Framework Application \\
3  & Designing for Humans: The Schema & Creator as Annotator \\
4  & The API Layer: Rust and Axum & Intervention Design \\
5  & Identity in the Loop: OIDC and Roles & Stakes Calibration \\
6  & Answer Validation: The Judgment Call & Threshold as Policy \\
7  & The Hint Engine: When to Help & Timing \\
8  & AI in the Loop: The Python Sidecar & Uncertainty Detection \\
9  & The Creator's Perspective & Human Contribution \\
10 & The Player's Perspective & Oracle as User \\
11 & Watching in Real Time & Feedback Integration \\
12 & QR, NFC, and Physical Anchors & Physical Loop Closure \\
13 & Testing HITL Systems & Feedback at Scale \\
14 & Deployment: Shipping the Loop & Real Stakes \\
15 & What the Game Taught Us & Meta-Reflection \\
\bottomrule
\end{tabularx}
\end{center}

\clearpage

% ── Milestones ────────────────────────────────────────────────────────────────
\section*{Milestone Cards}

Each card is one week's work. The \textbf{Build} column is concrete; the
\textbf{Reflect} prompt is not graded. The checkpoint is the deliverable you
should be able to demonstrate at the start of the following week.

\bigskip

% ----- Week 1 ----------------------------------------------------------------
\begin{milestone}{1}{Building HITL: Why a Game?}

\build{Nothing to code yet. Read \texttt{intro.md}. Run
\texttt{docker compose up -d --build}. Verify services start.
Browse the full directory tree and write a one-paragraph description of
each top-level folder in your own words.}

\medskip
\reflect{You have probably used a HITL system today. Where was the human?
Where was the loop? Did the loop actually close, or did it stop somewhere?}

\begin{hitlconcept}{The Loop Made Tangible}
  A treasure hunt is, structurally, a HITL system. By the end of this
  semester, you will not just know that — you will have felt it.
\end{hitlconcept}

\begin{checkpoint}
  Docker services healthy. Directory tree documented in your notebook.
  One-paragraph answer: what is the human's job in this system?
\end{checkpoint}

\end{milestone}

% ----- Week 2 ----------------------------------------------------------------
\begin{milestone}{2}{The Five Dimensions in Play}

\build{Read \texttt{chapters/02\_five\_dimensions.md} in full. Open
\texttt{backend-rust/src/services/clue\_validator.rs}. Annotate five
specific lines (add a comment) where a dimension of the framework is
embodied in the code. Do the same for one file in \texttt{backend-python/}.}

\medskip
\reflect{``The human is in the loop'' often means ``a human approves outputs
before they are acted on.'' Is approval the same as understanding?
When does it add value, and when is it theater?}

\begin{hitlconcept}{Framework Application}
  The Five Dimensions are not abstract. They live in specific files,
  specific parameters, specific defaults someone chose and stopped thinking
  about.
\end{hitlconcept}

\begin{checkpoint}
  Ten annotated lines (five Rust, five Python) with dimension labels and
  one sentence of reasoning each.
\end{checkpoint}

\end{milestone}

% ----- Week 3 ----------------------------------------------------------------
\begin{milestone}{3}{Designing for Humans: The Schema}

\build{Study \texttt{db/migrations/001\_initial.sql} and
\texttt{002\_external\_id.sql}. Run both against your local Postgres.
Sketch the entity-relationship diagram by hand. Identify every field where
a human made a normative decision (``what counts as a correct answer?'')
vs.\ a technical one (``which type stores this?''). Add one field you think
is missing.}

\medskip
\reflect{A database schema is a theory about the world. Every constraint is
a claim about what is always true. Which constraints here encode a claim
about human behavior?}

\begin{hitlconcept}{Creator as Annotator}
  The \texttt{answer\_tolerance} field is a threshold the creator sets.
  A default is a decision made once, propagated silently to everyone.
\end{hitlconcept}

\begin{hintbox}
  The \texttt{external\_id} column in migration 002 is for QR/NFC physical
  anchors. Don't try to understand it fully yet — it becomes clear in Week 12.
\end{hintbox}

\begin{checkpoint}
  ER diagram. Annotated list of normative vs.\ technical decisions.
  Proposed new field with justification.
\end{checkpoint}

\end{milestone}

% ----- Week 4 ----------------------------------------------------------------
\begin{milestone}{4}{The API Layer: Rust and Axum}

\build{Implement the \texttt{GET /hunts} and \texttt{POST /hunts} endpoints
in \texttt{src/routes/hunts.rs}. Add basic input validation. Write the
handler tests. Run \texttt{cargo test} until they pass. Then: try to make
a request that your validation should reject, and verify it is rejected.}

\medskip
\reflect{A type error at compile time is a design conversation. A type error
at runtime is an incident. Find one place where Rust's type system enforced
a HITL design decision before any player touched the system.}

\begin{hitlconcept}{Intervention Design}
  Every API contract is a decision about what the human can and cannot
  express. What can a creator \emph{not} say to this system, because the
  schema won't let them?
\end{hitlconcept}

\begin{checkpoint}
  \texttt{GET /hunts} and \texttt{POST /hunts} pass handler tests.
  One documented example of a rejected invalid input.
\end{checkpoint}

\end{milestone}

% ----- Week 5 ----------------------------------------------------------------
\begin{milestone}{5}{Identity in the Loop: OIDC and Roles}

\build{Configure Zitadel with creator and player dev accounts. Implement the auth middleware in
\texttt{src/middleware/auth.rs}. Verify that a player JWT cannot access a
creator-only endpoint. Document the RBAC rules in a table.}

\medskip
\reflect{Identity is not a property of a person. It is a claim made by a
system about a person, which other systems choose to trust. What happens
when the system's model of identity doesn't match the real-world situation?
(Two players sharing a phone. A creator who is also playing.)}

\begin{hitlconcept}{Stakes Calibration}
  Role separation is stakes calibration made structural. A player with
  creator access could invalidate everyone's data.
\end{hitlconcept}

\begin{checkpoint}
  Login works for both roles. Creator endpoint returns 403 for player JWT.
  RBAC table documented.
\end{checkpoint}

\end{milestone}

% ----- Week 6 ----------------------------------------------------------------
\begin{milestone}{6}{Answer Validation: The Judgment Call}

\build{Implement fuzzy matching in \texttt{src/services/clue\_validator.rs}
using the Levenshtein distance or Jaro-Winkler similarity. Make the
\texttt{answer\_tolerance} field live — different clues should behave
differently. Write a test matrix: ten (answer, correct\_answer) pairs,
three tolerance levels, expected result for each.}

\medskip
\reflect{Every threshold is an argument. The argument is usually implicit.
Write out the explicit argument your default threshold makes. Who benefits
from a strict threshold? Who is harmed?}

\begin{hitlconcept}{Threshold as Policy}
  The validation tolerance is not a technical parameter. It is a policy
  decision about how much linguistic variation you accept from players —
  which is a proxy for who you designed the game for.
\end{hitlconcept}

\begin{hintbox}
  The \texttt{strsim} crate has Jaro-Winkler implemented. Add it to
  \texttt{Cargo.toml} rather than implementing from scratch.
\end{hintbox}

\begin{checkpoint}
  Test matrix (30 cells) with all cases passing. Default tolerance
  documented and justified in one paragraph.
\end{checkpoint}

\end{milestone}

% ----- Week 7 ----------------------------------------------------------------
\begin{milestone}{7}{The Hint Engine: When to Help}

\build{Build the hint system: hints unlock after $n$ failed attempts
(configurable per clue). Implement three hint levels — vague, directional,
near-explicit. Track hint requests in the event log. Add a per-clue
\texttt{hint\_policy} field (none, after\_3, after\_5, always\_available).}

\medskip
\reflect{Help that arrives too early prevents learning. Help that arrives
too late causes frustration. The gap between them is the entire art of
pedagogy. Where does your default policy sit in that gap?
Who did you optimize for?}

\begin{hitlconcept}{Timing}
  The hint unlock timer is the most direct implementation of timing in
  this project. The threshold you choose is a claim about when frustration
  outweighs learning.
\end{hitlconcept}

\begin{checkpoint}
  Three-level hints working. Hint events logged with timestamps.
  Policy field configurable. Demo: one clue that unlocks hints at attempt 3,
  one that unlocks at attempt 5.
\end{checkpoint}

\end{milestone}

% ----- Week 8 ----------------------------------------------------------------
\begin{milestone}{8}{AI in the Loop: The Python Sidecar}

\build{Implement two services in \texttt{backend-python/}:
(1) \texttt{clue\_generator.py} — given a location description and theme,
generates clue text via the Anthropic API;
(2) \texttt{difficulty.py} — estimates the first-attempt success rate
for a clue using sentence embeddings and a trained regressor.
Connect both to the Rust API via HTTP. The sidecar must fail gracefully:
if it is down, the game still runs.}

\medskip
\reflect{An AI generates a clue. The creator approves it without reading
carefully. A player spends 40 minutes because the phrasing was subtly
ambiguous. Who is responsible?}

\begin{hitlconcept}{Uncertainty Detection}
  The difficulty estimator outputs a probability. What should the system
  do when it is uncertain about that estimate? When should it route the
  clue to the creator for review rather than publishing automatically?
\end{hitlconcept}

\begin{hintbox}
  Start with a lightweight similarity baseline for semantic matching and
  difficulty estimation. If you later need stronger embeddings, add them
  deliberately rather than making them a boot-time dependency.
\end{hintbox}

\begin{watchout}
  The Anthropic API has rate limits. Cache generated clues in Postgres,
  not in memory — the sidecar will restart.
\end{watchout}

\begin{checkpoint}
  \texttt{POST /clues/generate} returns a clue. \texttt{POST
  /clues/estimate-difficulty} returns a probability. Both endpoints return
  gracefully when sidecar is stopped.
\end{checkpoint}

\end{milestone}

% ----- Week 9 ----------------------------------------------------------------
\begin{milestone}{9}{The Creator's Perspective}

\build{Build the React creator dashboard in \texttt{web/}:
a hunt builder where the creator can create a hunt, add clues, set
tolerances, preview the clue sequence, and export a QR sheet (stub the QR
generation for now — it becomes real in Week 12). Integrate AI clue
generation so the creator can request suggestions and edit them.}

\medskip
\reflect{The annotator sees the world through the categories you gave them.
What categories does your creator dashboard give to the creator?
What can they not express because the form won't let them?}

\begin{hitlconcept}{Human Contribution}
  Every UI affordance is a constraint on what the human can contribute.
  The creator's dashboard shapes the quality of every player's experience
  before a single player has joined.
\end{hitlconcept}

\begin{checkpoint}
  Hunt buildable end-to-end via the creator UI. AI suggestions requestable
  and editable. Hunt visible in the database after save.
\end{checkpoint}

\end{milestone}

% ----- Week 10 ---------------------------------------------------------------
\begin{milestone}{10}{The Player's Perspective}

\build{Build the Flutter mobile app in \texttt{mobile/}:
the player scans a QR code (use the device camera), submits an answer,
sees the result, and can request a hint. The app should work offline for
content already downloaded; only answer submission requires connectivity.}

\medskip
\reflect{The oracle does not know they are an oracle. They are just trying
to win. Is there anything in your player UI that makes the data-collection
nature of the system visible? Should there be?}

\begin{hitlconcept}{Oracle as User}
  Every interaction the player makes is feedback. The quality of that
  feedback depends entirely on the quality of their experience. A confusing
  UI produces noisy labels.
\end{hitlconcept}

\begin{hintbox}
  Use \texttt{mobile\_scanner} for QR reading. Test on a physical device
  early — the simulator camera is unreliable.
\end{hintbox}

\begin{checkpoint}
  Player can scan a QR code, submit an answer, receive feedback, and
  request a hint. Full flow works on a physical device or emulator.
\end{checkpoint}

\end{milestone}

% ----- Week 11 ---------------------------------------------------------------
\begin{milestone}{11}{Watching in Real Time}

\build{Implement the WebSocket observer stream in
\texttt{src/routes/observer.rs}. Every significant event — clue scanned,
answer submitted (correct or not), hint requested, hunt completed —
broadcasts to connected observers. Build a minimal observer view in the
React dashboard: a live feed that shows events as they happen.}

\medskip
\reflect{You cannot understand a HITL system by reading its logs after the
fact. You understand it by watching it run. What does the live feed make
visible that the database would have hidden?}

\begin{hitlconcept}{Feedback Integration}
  The observer stream closes the feedback loop for the developer. Once you
  have real play data, re-train your difficulty estimator from Week 8 on
  actual attempt counts. How much did the synthetic training data mislead you?
\end{hitlconcept}

\begin{checkpoint}
  Live event stream works: three events visible in the observer view within
  10 seconds of a player action. Difficulty estimator re-trained on at least
  20 real play observations.
\end{checkpoint}

\end{milestone}

% ----- Week 12 ---------------------------------------------------------------
\begin{milestone}{12}{QR, NFC, and Physical Anchors}

\build{Implement QR sheet generation in the creator dashboard (use a
JavaScript QR library to generate printable A4 sheets with one QR per
clue location). If you have NFC hardware: implement NFC tag reading in
the Flutter app using \texttt{nfc\_manager}. Print and physically deploy
at least one three-clue hunt in a real space.}

\medskip
\reflect{The loop is not closed until a human has physically walked it.
What surprised you when you ran the hunt in a real space that no amount
of testing on a laptop would have revealed?}

\begin{hitlconcept}{Physical Loop Closure}
  A HITL system that has never touched a human is a simulation. Until
  someone walks the route with their phone, many bugs are invisible.
\end{hitlconcept}

\begin{hintbox}
  Use \texttt{qrcode} (npm) for QR generation. Print on A4, laminate if
  outdoors. NFC is optional — QR alone is sufficient for the checkpoint.
\end{hintbox}

\begin{checkpoint}
  Printable QR sheet generated for a three-clue hunt.
  Hunt successfully completed by at least one person who was not you.
  One documented observation from the live run.
\end{checkpoint}

\end{milestone}

% ----- Week 13 ---------------------------------------------------------------
\begin{milestone}{13}{Testing HITL Systems}

\build{Write integration tests that simulate the full loop:
(1) a creator creates a hunt with three clues;
(2) a simulated player makes three wrong attempts on clue 1, requests a hint,
then answers correctly;
(3) asserts the event log contains the expected sequence;
(4) asserts the hint was unlocked at the right attempt count.
Also: write one test that deliberately injects a flawed clue and asserts
the system handles it gracefully.}

\medskip
\reflect{A test that passes tells you the code does what you expected.
It tells you nothing about whether what you expected was right.
What would a test of your \emph{design assumptions} look like?}

\begin{hitlconcept}{Feedback at Scale}
  The integration test is a synthetic human. It is faster and cheaper than
  a real player, but it only tests the behaviors you thought to specify.
  Where is your test suite blind?
\end{hitlconcept}

\begin{checkpoint}
  Integration test suite covers the full three-clue loop with wrong
  attempts, hint unlock, and correct answer. All tests pass in CI
  (\texttt{cargo test}).
\end{checkpoint}

\end{milestone}

% ----- Week 14 ---------------------------------------------------------------
\begin{milestone}{14}{Deployment: Shipping the Loop}

\build{Deploy the full stack to a machine that is not your laptop —
a VPS, a local server, or a Raspberry Pi on a local network all count.
Participants should be able to join using their own phone without any
special setup. Run a live hunt with at least three participants you did
not brief in advance.}

\medskip
\reflect{A system that works on your laptop is a prototype. A system that
works for strangers is a product. What was the first thing that broke when
strangers used it?}

\begin{hitlconcept}{Real Stakes}
  Deployment changes the cost model. Errors are no longer free. A
  misconfigured threshold, a missing hint, a confusing UI — all of these
  now have real consequences for real people's experience.
\end{hitlconcept}

\begin{watchout}
  If your VPS has limited RAM, the Rust binary + Postgres + Python sidecar
  together use \textasciitilde 800MB. A 1GB instance will OOM under load.
  Use a 2GB minimum or run Postgres on a separate container.
\end{watchout}

\begin{checkpoint}
  System deployed and reachable from a device not on your local network.
  Live hunt completed by at least three external participants.
  Notes from what broke, what surprised you, what you changed under pressure.
\end{checkpoint}

\end{milestone}

% ----- Week 15 ---------------------------------------------------------------
\begin{milestone}{15}{What the Game Taught Us}

\build{Write a retrospective analysis (1,500--2,000 words). Structure it
around the Five Dimensions: for each dimension, describe one decision you
made in the code, whether it worked, and what you would change. End with
one paragraph on what the game made visible that the textbook could not.}

\medskip
\reflect{You do not understand a system by reading about it.
You understand it by breaking it.
What did you understand, after Week 14, that you did not understand
after Week 1?}

\begin{hitlconcept}{Meta-Reflection}
  The retrospective is not a formality. It is the closing of the loop
  you have been building for fifteen weeks: human feedback on the system
  that was designed to generate human feedback.
\end{hitlconcept}

\begin{checkpoint}
  Retrospective submitted (1,500--2,000 words).
  Five Dimensions each addressed with a concrete code example.
  One ``this surprised me'' paragraph grounded in the live deployment.
\end{checkpoint}

\end{milestone}

\clearpage

% ── Final Integration Checklist ───────────────────────────────────────────────
\section*{Final Integration Checklist}

Before your Week 15 retrospective, verify the following end-to-end:

\begin{tcolorbox}[
  enhanced,
  colframe=SpringLeaf!50,
  colback=white,
  top=8pt, bottom=8pt, left=12pt, right=12pt,
  arc=4pt,
]
\renewcommand{\arraystretch}{1.4}
\begin{tabularx}{\linewidth}{@{}>{\footnotesize}clX@{}}
\toprule
& & \textbf{Capability} \\
\midrule
\faSquare[regular] & & Creator can log in and build a hunt via the React dashboard \\
\faSquare[regular] & & Creator can request AI clue suggestions and edit them \\
\faSquare[regular] & & Creator can export a printable QR sheet \\
\faSquare[regular] & & Player can scan a QR code on a physical device \\
\faSquare[regular] & & Player can submit an answer and receive correct/incorrect feedback \\
\faSquare[regular] & & Answer validation respects the per-clue tolerance setting \\
\faSquare[regular] & & Hints unlock after the configured number of failed attempts \\
\faSquare[regular] & & Observer sees live event stream during active play \\
\faSquare[regular] & & AI sidecar failure does not crash the core game \\
\faSquare[regular] & & Difficulty estimates update when re-trained on real play data \\
\faSquare[regular] & & Integration test suite passes on a clean checkout \\
\faSquare[regular] & & System survives a 30-minute live hunt with $\geq$ 3 participants \\
\faSquare[regular] & & Retrospective maps each deliverable to a Five Dimensions concept \\
\bottomrule
\end{tabularx}
\end{tcolorbox}

\bigskip

% ── The Five Dimensions Reference ─────────────────────────────────────────────
\section*{Reference: The Five Dimensions}

\begin{tcolorbox}[
  enhanced,
  colframe=ConceptBlue!40,
  colback=ConceptIce!60,
  top=8pt, bottom=8pt, left=12pt, right=12pt,
  arc=4pt,
]
\renewcommand{\arraystretch}{1.5}
\begin{tabularx}{\linewidth}{@{}>{\bfseries\color{ConceptBlue}\small}lX@{}}
\toprule
Dimension & Key Question \\
\midrule
Uncertainty Detection  & Can the system recognize when it is unsure? \\
Intervention Design    & How does it ask for help — or choose not to? \\
Timing                 & When does it intervene — too early, too late, or just right? \\
Stakes Calibration     & Does it understand the cost of being wrong? \\
Feedback Integration   & Does it learn from the responses it receives? \\
\bottomrule
\end{tabularx}
\end{tcolorbox}

\bigskip

\begin{center}
  {\color{RuleGray}\rule{0.6\linewidth}{0.4pt}}

  \vspace{0.4cm}

  {\small\sffamily\color{gray}
    Human in the Loop\enspace —\enspace Spring Edition\enspace —\enspace
    \textit{Build the Loop}}

  \vspace{0.2cm}

  {\footnotesize\color{gray}
    laztoum.com\enspace·\enspace what-if.io\enspace·\enspace limen-os.io\enspace·\enspace laztoum@protonmail.com}
\end{center}

\end{document}
